% Complete Chapter 6 introduction, original PDF 73 / printed 63.
% block: exterior-calculus-chapter-heading page:73 type:chapter
\clearpage\setcounter{section}{5}
\section{\BT{Basic notions of exterior calculus}{外微积分的基本概念}}\label{chap:exterior-calculus}

% block: exterior-calculus-chapter-introduction page:73 type:ordinary-paragraph
\BT{We have thus far studied Hilbert complexes in the abstract and mostly illustrated the concepts via the de Rham complex on a domain in three dimensions. In that example the complex has four spaces and three differential operators, which are the vector calculus operators gradient, curl, and divergence, whose properties we studied individually, employing very analogous arguments for each. In this second part of the book, we will focus on the de Rham complex in its most natural setting, on an $n$-dimensional manifold. The complex then has $n+1$ spaces connected by $n$ distinct operators. Using the language of exterior calculus we can introduce all these operators in a uniform way, obtaining a theory which is at once more elegant, more general, and more insightful. In this chapter, we present the relevant exterior calculus, culminating in the $L^2$ de Rham complex on a Riemannian manifold. Then, in the following chapter, we turn to the discretization of the corresponding Hodge Laplace boundary value problems using finite elements.}{到目前为止, 我们抽象地研究了 Hilbert 复形, 并主要用三维区域上的 de Rham 复形说明这些概念. 在该例中, 复形有四个空间和三个微分算子, 即向量微积分中的梯度, 旋度和散度算子; 我们逐个研究了它们的性质, 对每个算子采用了十分相似的论证. 在本书的第二部分, 将着重研究 de Rham 复形最自然的情形, 即 $n$ 维流形上的情形. 此时复形有 $n+1$ 个空间, 由 $n$ 个不同的算子连接. 使用外微积分的语言, 可以统一地引入所有这些算子, 从而得到更优美, 更一般且更有洞察力的理论. 本章介绍相关的外微积分, 最后讨论 Riemannian 流形上的 $L^2$ de Rham 复形. 随后在下一章, 转向用有限元离散化相应的 Hodge Laplace 边值问题.}
